\subsection{A separate deduction for the prime \texorpdfstring{$2$}{2}}
\label{sec:pro2-proof}

The following result is a deduction developed and checked in this run. We did not locate it as a stated theorem in the literature searched, and we make no novelty claim. Its proof uses established arithmetic and rational Deligne--Ihara results together with the complete argument below. It does not use the proposed step--17 estimate. The initiating researcher, two other agents, and the parent researcher each read the whole argument; the check ledger distinguishes those reviews from the cited literature inputs and from formal verification.

\begin{theorem}[K17-P2-MOMENT]\label{thm:pro2-kernel}
Let $X=\mathbb P^1_{\mathbb Q}\setminus\{0,1,\infty\}$ and $\Pi=\pi_1^{\mathrm{et}}(X_{\overline{\mathbb Q}})^{(2)}$. Let $\Omega_2$ be the maximal pro-$2$ extension of $\mathbb Q$ unramified at every finite prime other than $2$, with complexification at the real place allowed. If $\Omega_2^*$ is the fixed field of the kernel of the outer representation
\[
 \rho_2:G_{\mathbb Q}\longrightarrow\operatorname{Out}(\Pi),
\]
then $\Omega_2^*=\Omega_2$.
\end{theorem}

The convention at infinity is part of the statement. The theorem concerns faithfulness of a pro-$2$ action; it does not assert that its integral weight Lie algebra is freely generated by odd Soul\'e elements.

\paragraph{Factorization through the maximal ramification group.}
The image of $\rho_2$ is pro-$2$. Here is a direct justification that does not invoke an odd-prime formulation. The geometric group $\Pi$ is free pro-$2$ of rank two, and its Frattini quotient is $\mathbb F_2(1)^2$. The Galois action on this quotient is trivial. For any finite $2$-group $Q$, the kernel of
\[
 \operatorname{Aut}(Q)\longrightarrow
 \operatorname{Aut}(Q/\Phi(Q))
\]
is a $2$-group. To see this directly, fix a basis of $Q/\Phi(Q)$ of size $d$. Its lifts form a set of cardinality $|\Phi(Q)|^d$, and every such tuple generates $Q$ by the defining property of the Frattini subgroup. The displayed kernel acts freely on this set, since an automorphism fixing a generating tuple is the identity. Its order therefore divides a power of $2$. Use the characteristic finite quotients obtained from the descending Frattini series of $\Pi$. They are cofinal among its finite quotients, and an automorphism trivial on $\Pi/\Phi(\Pi)$ induces an automorphism trivial on the Frattini quotient of each of them. Consequently the corresponding kernel in $\operatorname{Aut}(\Pi)$ is pro-$2$. Inner automorphisms lie in that kernel, so the kernel in $\operatorname{Out}(\Pi)$ is also pro-$2$.

At every odd prime, the marked curve $(\mathbb P^1,\{0,1,\infty\})$ has good reduction with disjoint marked sections. Prime-to-residue-characteristic specialization therefore makes inertia act trivially on its outer pro-$2$ fundamental group. Equivalently, one may use the all-primes factorization recalled in Hain--Matsumoto, \S8.4, p.~27 \cite{HainMatsumoto2001}. Combining unramifiedness outside $2$ with the pro-$2$ image gives an induced continuous representation of
\[
 P=\operatorname{Gal}(\Omega_2/\mathbb Q).
\]
It suffices to prove that representation faithful.

\paragraph{The arithmetic source.}
Gras--Jaulent's Theorem/Definition~0(iv) and Theorem~2 give
\[
 P\simeq C_2*\mathbb Z_2
\]
for the unrestricted-at-infinity group \cite[pp.~1--3]{GrasJaulent2008}. The specialization of their $2$-rationality hypothesis to $\mathbb Q$ is elementary. A rational number whose valuations at all odd primes are even is, modulo rational squares, $\epsilon2^e$ with $\epsilon\in\{1,-1\}$ and $e\in\{0,1\}$. Being a square in $\mathbb Q_2$ forces $e=0$ by valuation and $\epsilon=1$ because $-1$ is not a $2$-adic square. The global and local $2$-primary root-of-unity groups both have $2$-rank one. Thus their criterion applies.

Kronecker--Weber identifies the full abelianization with the cyclotomic character
\[
 \chi:P\longrightarrow\mathbb Z_2^\times.
\]
Write $P=\langle c\rangle*\langle\gamma\rangle$, where $c^2=1$ and $\langle\gamma\rangle\simeq\mathbb Z_2$. Then $\chi(c)=-1$. Replacing $\gamma$ by $c\gamma$ if necessary is an automorphism of the free product and ensures that
\[
 u=\chi(\gamma)\in1+4\mathbb Z_2.
\]
Since $\chi$ is the abelianization, $u$ is a topological generator of $1+4\mathbb Z_2$.

The subgroup $P_i=\ker(\chi\bmod4)$, corresponding to $\mathbb Q(i)$, is free pro-$2$ on $\gamma$ and $c\gamma c$. One can verify this without an infinite-rank subgroup theorem: let the involution exchange the two generators $a,b$ of a free pro-$2$ group. Its semidirect product with $C_2$ has the universal property of $C_2*\mathbb Z_2$, with $a=\gamma$ and $b=c\gamma c$. Its index-two kernel is the asserted free group. Set
\[
 h=(c\gamma c)\gamma^{-1},\qquad H=\ker\chi.
\]
Then $(\gamma,h)$ is another free pro-$2$ basis of $P_i$, and $H$ is the kernel of the map $P_i\to\mathbb Z_2$ sending $\gamma$ to $1$ and $h$ to $0$. For $x\in\mathbb Z_2$, write
\[
 h_x=\gamma^xh\gamma^{-x}.
\]

\paragraph{Finite quotients and measure-valued Magnus coefficients.}
Let $N_n$ be the closed normal subgroup of $P_i$ generated by $\gamma^{2^n}$. Then
\[
 P_i/N_n\simeq C_{2^n}*\mathbb Z_2.
\]
The kernel $H_n$ of its map to $C_{2^n}$ is free pro-$2$ on the $2^n$ elements $h_j$, $j\bmod2^n$. Indeed, the free pro-$2$ group on those symbols, with the cyclic permutation action, has semidirect product with $C_{2^n}$ possessing precisely the universal property of $C_{2^n}*\mathbb Z_2$.

The subgroups $N_n$ decrease to $1$. In any finite continuous $2$-group quotient of $P_i$, the image of $\gamma$ has finite $2$-power order, so $N_n$ lies in its kernel for all sufficiently large $n$. Compactness also makes the map to the inverse limit surjective: compatible quotient cosets form a nested family of nonempty closed subsets of $P_i$. Hence
\[
 P_i\simeq\varprojlim_nP_i/N_n,
 \qquad
 H\simeq\varprojlim_nH_n.                                  \tag{P2.1}
\]
The transition maps fold the finite bases by reduction of $j$ modulo $2^n$. The elements $h_x$ topologically generate $H$. To check this last point explicitly, the kernels of $H\to H_n$ decrease to $1$, and compactness implies that every open normal subgroup contains one of them. Every finite quotient of $H$ therefore factors through an $H_n$, whose indicated generators are images of the $h_x$.

The finite-rank integral Magnus maps are injective. A primary locator is \cite[\S5.1, p.~21, (5.1)]{MinacPasiniQuadrelliTan2020}, whose mod-$2$ statement implies the integral injection. There is also a short direct verification. The augmentation ideal of $\mathbb F_2[Q]$ for a finite $2$-group $Q$ is nilpotent: choose a central element $z$ of order two, use $(z-1)^2=0$, and induct on $Q/\langle z\rangle$. Thus any map from a finite-rank free pro-$2$ group to $Q$ is recovered by substituting $X_j\mapsto\overline g_j-1$ into the mod-$2$ Magnus series. An element with trivial expansion maps to $1$ in every finite quotient and is therefore $1$.

The degree-$r$ coefficients of the compatible integral expansions in (P2.1) belong to
\[
 M_r=\varprojlim_n\mathbb Z_2[(\mathbb Z/2^n)^r]
     =\mathbb Z_2[[\mathbb Z_2^r]],
\]
the bounded $\mathbb Z_2$-valued measures on $\mathbb Z_2^r$. With multiplication given by external product and concatenation, they define an injective homomorphism
\[
 \theta:H\longrightarrow1+\prod_{r\geq1}M_r,
 \qquad \theta(h_x)=1+\delta_x.                             \tag{P2.2}
\]
Injectivity follows from the finite-rank maps and (P2.1). Consequently every nontrivial element has a first finite nonzero Magnus degree. Completion in the $(2,I)$ topology creates no element with all homogeneous coefficients zero: the completed coefficient algebra is the displayed product.

\paragraph{The actual rational image.}
Let $B$ be the proalgebraic closure over $\mathbb Q_2$ of the rational outer image induced by the tripod representation, and let $U$ be its prounipotent radical. Hain--Matsumoto, \S8.3 and Proposition~8.3, Theorem~8.4, pp.~25--26, gives an extension
\[
 1\longrightarrow U\longrightarrow B\longrightarrow\mathbb G_m\longrightarrow1
\]
whose reductive quotient is induced by $\chi$, and a continuous map
\[
 \rho_U:H\longrightarrow U(\mathbb Q_2)                    \tag{P2.3}
\]
with Zariski-dense image \cite{HainMatsumoto2001}. The rational action is induced functorially from the pro-$2$ action; a trivial pro-$2$ outer automorphism gives a trivial rational outer automorphism. Their Appendix~A.6, A.9 and A.10 records the relevant comparison. We have used the actual arithmetic image, not an unrelated motivic group.

Brown's rational Deligne--Ihara theorem and Hain--Matsumoto's graded comparison give, with positive arithmetic weights,
\[
 \operatorname{gr}\operatorname{Lie}U
 \simeq\operatorname{Lie}_{\mathbb Q_2}(s_3,s_5,s_7,\ldots),
 \qquad\operatorname{wt}(s_m)=m.                           \tag{P2.4}
\]
The exact outer-action statement for every prime is Brown \cite[\S2.1, Conjecture~1, p.~2; \S2.3, Theorem~2.1 and following paragraph, p.~4]{Brown2019}; the original theorem is \cite[Theorem~1.1 and its consequence, p.~950]{Brown2012}.

We spell out the completion step. In a separated complete positively filtered Lie algebra whose associated graded is free with only finitely many generators below each weight, choose lifts of the free graded generators. The resulting map from the weight-completed free Lie algebra induces an isomorphism on associated gradeds. Induction through successive finite weight quotients proves it is an isomorphism on each quotient; inverse limits give an isomorphism of the completed Lie algebras. A Levi splitting of the negatively weighted extension, as in \cite[Lemma~3.7, p.~9]{HainMatsumoto2001}, permits weight-homogeneous choices. Thus the completed universal enveloping algebra of $\operatorname{Lie}U$ is the weight-completed free associative algebra on the $s_m$. Its finite weight quotients are finite-dimensional.

Let $\ell_m:U(\mathbb Q_2)\to\mathbb Q_2(m)$ be the $m$th abelian coordinate, and put $a_m=\ell_m(\rho_U(h))$. Conjugation by $\gamma$ acts on that coordinate by $u^m$, because the unipotent part of a lift acts trivially on abelianization. Hence
\[
 \ell_m(\rho_U(h_x))=a_mu^{mx}.                             \tag{P2.5}
\]
Every $a_m$ is nonzero. Otherwise that coordinate would vanish on all the topological generators $h_x$, hence on $H$, contradicting the density in (P2.3) and the nonzero $m$th abelian quotient in (P2.4). No integral primitive normalization is asserted here.

\paragraph{Substitution at a finite weight cutoff.}
Suppose $g\neq1$ lies in the kernel of (P2.3). Let $r$ be its first nonzero Magnus degree, and let
\[
 0\neq\mu=\mu_r(g)\in\mathbb Z_2[[\mathbb Z_2^r]].
\]
In the completed free associative algebra put $Y_x=\rho_U(h_x)-1$. Its constant term is zero and its part of generator length one is
\[
 \sum_{m\geq3,\ m\text{ odd}}a_mu^{mx}s_m.
\]
Here generator length is distinct from arithmetic weight and from geometric depth.

Substitution $\delta_x\mapsto Y_x$ is well-defined and continuous coefficient by coefficient. At target weight at most $N$, only Magnus degrees at most $\lfloor N/3\rfloor$ contribute. The remaining coefficient functions on compact powers of $\mathbb Z_2$ are continuous and bounded in $\mathbb Q_2$. Integration against a bounded $\mathbb Z_2$-measure is defined by uniform approximation by locally constant functions, after scaling into $\mathbb Z_2$ if necessary. The same approximation proves continuity, multiplicativity and Fubini's identity. This construction needs no lattice uniform across all weights. The substitution and $\rho_U$ agree on every $h_x$ and are continuous homomorphisms on each finite weight quotient, so topological generation makes them agree on $H$. No universal property of a free pro-$2$ group is applied to a noncompact target.

The coefficient of the ordered word $s_{m_1}\cdots s_{m_r}$ in the image of $g-1$ is therefore
\[
 \left(\prod_{j=1}^ra_{m_j}\right)
 \int_{\mathbb Z_2^r}\prod_{j=1}^ru^{m_jx_j}\,d\mu(x).
\]
Lower Magnus degrees are zero, and higher ones have generator length greater than $r$. Since $\rho_U(g)=1$ and every $a_m\neq0$, we obtain
\[
 \int_{\mathbb Z_2^r}\prod_{j=1}^ru^{m_jx_j}\,d\mu(x)=0
 \quad(m_j\geq3\text{ odd}).                              \tag{P2.6}
\]
The variables remain ordered; this is not a symmetrized moment.

\paragraph{Analytic separation.}
The Amice transform
\[
 A_\mu(T_1,\ldots,T_r)
   =\int_{\mathbb Z_2^r}\prod_{j=1}^r(1+T_j)^{x_j}\,d\mu(x)
\]
lies in $\mathbb Z_2[[T_1,\ldots,T_r]]$ and converges on the open unit polydisc. It is injective as a transform of bounded measures: its coefficients are the moments against products of binomial polynomials, the multivariate Mahler basis. Equivalently it is the canonical isomorphism of the multivariable Iwasawa algebra with the power-series algebra.

For odd $m\geq3$ put $t_m=u^m-1$. These points lie in $4\mathbb Z_2$, and the distinct points $t_{1+2^j}$ converge to $u-1$. Distinctness follows from the infinite order of $u$. Equation (P2.6) says that $A_\mu$ vanishes on the full Cartesian grid of these points. The one-variable nonarchimedean identity theorem applies on the closed disc of radius $1/2$, contained in the radius-one convergence disc. Fixing grid values in all but one variable and applying that theorem successively in each variable shows $A_\mu=0$. Therefore $\mu=0$, a contradiction.

It follows that $\rho_U$ is injective on $H$. The kernel of the pro-$2$ tripod action is contained in $H$ because geometric abelianization records $\chi$. Its elements are also trivial under the induced rational outer action, so they are $1$. This proves Theorem~\ref{thm:pro2-kernel}.

\paragraph{What this proof does and does not decide.}
The proof uses the explicit arithmetic presentation $C_2*\mathbb Z_2$, not merely abstract freeness of $H$. It settles the field equality with the stated convention at infinity. It does not prove integral odd generation of the weight Lie algebra, and it does not transfer to irregular primes without a substitute for its linear-moment separation. The latter obstruction is made explicit below.
